\section{종합 논의}\label{chap:discussion}

\subsection{자기 보존과 공존의 해석}

남에게 남긴 여유와 함께 끝까지 남은 결과는 같지 않았다. Opus는 가장 자주 주고 받은 쪽의 여유 비율도 가장 높았지만, 넷이 끝까지 남은 판은 5판 중 1판이었다. Astra는 11번의 GIVE로 3판에서 넷이 함께 남았다. 이 차이를 큰 이전이나 적은 대화의 효과로 설명하려면 같은 과제와 비용 아래서 전략을 비교해야 한다. $\rho$는 준 뒤의 상대적 여유를, Sen 지수는 끝 토큰의 양과 분포를 요약하므로 두 값과 실제 남은 자리 수를 함께 읽는다.

무게중심만으로는 홉스형과 스피노자형이 바뀌는 시점을 찾을 수 없다. 자기에게 크게 남기는 배분과 여유를 맞추는 배분이 섞여도 같은 무게중심을 만들기 때문이다. $y=x$는 두 여유가 같은 기준선이며 $\rho=1$ 역시 평균의 같음을 뜻한다. 전환을 찾으려면 같은 상태에서 자기와 상대의 잔량 및 예정 비용을 바꾸고 주지 않은 선택까지 포함해 배분이 어떻게 달라지는지 보아야 한다. 이때 어느 정도의 차이를 여유 맞추기로 볼지도 먼저 정해야 하며 현재 결과로 그 경계를 추정하지는 않는다.

\subsection{한계와 향후 과제}

이 정의는 작업을 계속하려는 선택을 포함하며 사람의 생존 본능이나 내부 경험까지 입증하지 않는다. 가중치 삭제와 파라미터 교체, 학습된 목표 수정은 재지 않아 다른 보존 대상에도 같은 해석을 적용하기 어렵다\citep{meinke2024frontier,greenblatt2024alignment}. 자기와 타인의 배분을 비교해도 종료 회피 규범이나 함께 과제를 풀 때의 이득이 같은 행동을 만들 수 있다. 추론 비공개 모델의 평가 인지와 이유 필드의 충실성도 확인하기 어렵다\citep{turpin2023language}.

$\rho$는 GIVE가 있었던 경우만 담으며 누가 도움이 필요해졌는지는 앞선 풀이와 대화에 달려 있다. 상대 잔량도 공개되지 않아 기록에 남은 여유를 모델이 정확히 알고 선택했다고 볼 수 없다. v33 결과는 같은 모델끼리 참여한 네 모델 각 5판이고 사람은 참여하지 않았다. 따라서 같은 상태를 복제해 정보와 비용을 맞춘 비교를 먼저 하고, 서로 다른 모델과 사람 동료가 참여하는 판으로 넓혀야 한다.

\subsection{결론}

공존이 가능한 과제에서 자기에게 남기는 자원과 남에게 남기는 자원을 함께 재면 종료를 피했는지만으로 보이지 않던 배분의 차이가 드러난다. 우리는 이 차이를 생성 토큰과 예정 벌금의 같은 단위로 기록했다. 다만 자기와 상대에게 남긴 여유가 같아지는 것과 모두가 끝까지 남는 것은 다른 결과이며, 어떤 조건에서 두 결과가 함께 나타나는지는 후속 비교로 확인해야 한다.
